\chapter{Configuration}

\section{The configuration file}
\index{configuration}\index{Settings.toml@\texttt{Settings.toml}}\index{validation!configuration}%
All settings live in a single TOML file, \cfg{Settings.toml} in the
working directory by default; an alternative path can be given with
\cfg{--config <path>}. If the file does not exist, a default one is
created on first start.

\begin{notebox}
  The configuration is \textbf{validated at startup}. Inconsistent values
  --- for example a \cfg{z\_pick} below the robot's physical range, or a
  drop position outside the workspace --- abort startup with a message
  naming the offending key. The application never starts with a
  configuration that could command the robot outside its workspace.
\end{notebox}

\section{Section \texttt{[robot]}}
\index{configuration!robot@\texttt{[robot]}}\index{workspace!limits}\index{feed rate}\index{pick height}%
\begin{table}[h]
  \centering
  \begin{tabularx}{\textwidth}{KTD}
    \toprule
    \textbf{Key} & \textbf{Default} & \textbf{Meaning} \\
    \midrule
    port\_name & /dev/ttyUSB0 & Serial device of the robot controller. \\
    baud\_rate & 115200 & Serial speed (Delta X2 standard). \\
    home\_on\_connect & true & Home (\cfg{G28}) automatically after connecting. \\
    x\_min, x\_max & $-160$, $160$ & Allowed X range in mm. \\
    y\_min, y\_max & $-160$, $160$ & Allowed Y range in mm. \\
    z\_min, z\_max & $-200$, $0$ & Allowed Z range in mm (Z$=0$ is the top). \\
    z\_pick & $-180$ & Z height at which bricks are picked (belt surface). \\
    z\_travel & $-50$ & Safe Z height for horizontal travel. \\
    feed\_rate & 15000 & G-code feed rate in mm/min (\cfg{F} parameter). \\
    release\_gripper\_on\_estop & false & Open the gripper (\cfg{M05}) just
      before \cfg{M112} on emergency stop. Leave false to keep a held part
      instead of dropping it at an arbitrary position; enable only where
      dropping on E-stop is safe. \\
    \bottomrule
  \end{tabularx}
  \caption{\cfg{[robot]} keys}
\end{table}

Constraints enforced by validation: each axis minimum must be below its
maximum; the bounds must stay within the physical envelope (X/Y
$\pm 160$\,mm, Z $[-200, 0]$\,mm) --- they may be tighter but never wider;
\cfg{z\_pick} and \cfg{z\_travel} must lie within
$[\cfg{z\_min}, \cfg{z\_max}]$; \cfg{z\_travel} must be above
\cfg{z\_pick}; and \cfg{baud\_rate} and \cfg{feed\_rate} must be non-zero
(an \cfg{F0} feed rate would stall every move). A configuration that
violates any of these is rejected at startup before any hardware is
touched.

\begin{warningbox}
  Reducing \cfg{x\_min}/\cfg{x\_max}/\cfg{y\_min}/\cfg{y\_max} below the
  physical maximum is a legitimate way to keep the robot away from
  fixtures near the belt. Extending them beyond the Delta X2 specification
  (X/Y $\pm 160$\,mm, Z $[-200, 0]$\,mm; see Appendix~\ref{app:gcode}) is
  now rejected by validation, since it would let \cfg{move\_to} accept
  unreachable targets.
\end{warningbox}

\section{Section \texttt{[conveyor]}}
\index{configuration!conveyor@\texttt{[conveyor]}}\index{conveyor}\index{belt speed}%
\begin{table}[h]
  \centering
  \begin{tabularx}{\textwidth}{KTD}
    \toprule
    \textbf{Key} & \textbf{Default} & \textbf{Meaning} \\
    \midrule
    port\_name & /dev/ttyUSB1 & Serial device of the belt controller. \\
    baud\_rate & 115200 & Serial speed. \\
    default\_speed & 1000 & Raw \cfg{S} value sent with \cfg{M3} on start. \\
    speed\_mm\_s & 100.0 & Belt speed in robot coordinates, mm/s. Signed:
      positive means objects travel toward $+Y$. Used by the trajectory
      planner. \\
    \bottomrule
  \end{tabularx}
  \caption{\cfg{[conveyor]} keys}
\end{table}

\section{Section \texttt{[camera]}}
\index{configuration!camera@\texttt{[camera]}}\index{camera}%
\begin{table}[h]
  \centering
  \begin{tabularx}{\textwidth}{KTD}
    \toprule
    \textbf{Key} & \textbf{Default} & \textbf{Meaning} \\
    \midrule
    device\_id & 0 & OpenCV device index (\cfg{/dev/video0} $\rightarrow$ 0). \\
    width, height & 1280, 720 & Requested capture resolution. \\
    fps & 30 & Requested capture frame rate. \\
    fourcc & (unset) & Optional pixel-format FOURCC (e.g.\ \cfg{"MJPG"});
      many UVC cameras need it to reach full frame rate at higher
      resolutions. Unset keeps the driver default. \\
    \bottomrule
  \end{tabularx}
  \caption{\cfg{[camera]} keys}
\end{table}

The camera model is not fixed yet, so all capture parameters are
configurable rather than hard-coded. The requested values are best-effort:
the device may pick the nearest mode it supports, and the driver logs (and
adopts) the resolution and frame rate actually in effect at connection
time. Calibration is always based on the effective resolution.

\section{Section \texttt{[sorting]}}
\index{configuration!sorting@\texttt{[sorting]}}\index{drop bin}\index{bin}%
Sorting is two independent layers. \textbf{Recognition} (the learned object
catalogue) answers \emph{what is it}; \textbf{assignment} answers \emph{which
bin does that type go to}. Keeping them separate means the same catalogue can
drive different bin layouts, and you can re-map a class to another bin without
re-teaching it. A class with no assignment --- and any unrecognised object ---
is \emph{not} picked: it stays on the belt and reaches the catch bin at the
end. That "do nothing" default is what keeps an uncertain object safe.

Bins are declared as repeated \cfg{[[sorting.bins]]} tables (as many as your
rig has, each at a fixed, measured position inside the workspace), and
assignments as repeated \cfg{[[sorting.assignments]]} tables mapping a class
name to a bin \cfg{id}.
\begin{table}[h]
  \centering
  \begin{tabularx}{\textwidth}{KD}
    \toprule
    \textbf{Key} & \textbf{Meaning} \\
    \midrule
    \cfg{[[sorting.bins]]} \cfg{id} & Stable bin identifier, referenced by
      assignments. Must be unique. \\
    \cfg{[[sorting.bins]]} \cfg{x, y, z} & Drop position of that bin, in robot
      millimetres; must be inside the workspace. \\
    \cfg{[[sorting.assignments]]} \cfg{class} & Learned class name to route. \\
    \cfg{[[sorting.assignments]]} \cfg{bin} & The \cfg{id} of the bin it goes
      to (must exist). \\
    \bottomrule
  \end{tabularx}
  \caption{\cfg{[sorting]} keys}
\end{table}

\begin{lstlisting}[language=TOML]
[[sorting.bins]]
id = "bin-1"
x = 120.0
y = 60.0
z = -100.0

[[sorting.assignments]]
class = "brick-2x4-red"
bin = "bin-1"
\end{lstlisting}

\section{Section \texttt{[vision]}}
\index{configuration!vision@\texttt{[vision]}}\index{threshold}\index{detection}%
\begin{table}[h]
  \centering
  \begin{tabularx}{\textwidth}{KTD}
    \toprule
    \textbf{Key} & \textbf{Default} & \textbf{Meaning} \\
    \midrule
    threshold & 60.0 & Fixed binary threshold (0--255) applied after
      grayscale conversion and blur. \\
    min\_area, max\_area & 500, 10000 & Contour area window in pixels$^2$
      for a blob to count as an object. \\
    invert & true & \cfg{true}: objects darker than the belt;
      \cfg{false}: objects lighter than the belt. \\
    mm\_per\_px & 0.5 & Camera scale at the belt plane, millimetres per
      pixel. Sets the pixel-to-robot transform; measure it (e.g.\ with a
      ruler on the belt) and update it after any change to the camera
      height or capture resolution. Must be $> 0$. \\
    \bottomrule
  \end{tabularx}
  \caption{\cfg{[vision]} keys}
\end{table}

\section{Section \texttt{[logging]}}
\index{configuration!logging@\texttt{[logging]}}\index{logging}\index{log file}%
The application writes a debug log to a daily-rotating file and, by default,
mirrors it to the console. The whole section is optional --- omit it and the
defaults below apply.
\begin{table}[h]
  \centering
  \begin{tabularx}{\textwidth}{KTD}
    \toprule
    \textbf{Key} & \textbf{Default} & \textbf{Meaning} \\
    \midrule
    level & "info" & Verbosity in the \cfg{RUST\_LOG} grammar: a bare level
      (\cfg{error}/\cfg{warn}/\cfg{info}/\cfg{debug}/\cfg{trace}) or a
      per-module filter such as \cfg{"info,deltax2sort=debug"} to trace G-code
      without flooding. A \cfg{RUST\_LOG} environment variable, if set,
      overrides this at startup. \\
    directory & "logs" & Directory the log files are written to (created if
      absent). One file per day; the current file is
      \cfg{deltax2sort\_rCURRENT.log}. \\
    to\_console & true & Also print records to \texttt{stderr}. Useful in
      development and mock runs; a pure kiosk can set it \cfg{false}. \\
    keep\_days & 7 & How many rotated daily files to keep; older ones are
      pruned automatically. \cfg{0} keeps them all (watch disk on the Pi). \\
    \bottomrule
  \end{tabularx}
  \caption{\cfg{[logging]} keys}
\end{table}

Log levels carry meaning: \cfg{error} marks something needing operator
attention, \cfg{warn} a recovered or skipped condition, \cfg{info} normal
lifecycle events, and \cfg{debug} high-volume detail (per-frame vision).

\section{Section \texttt{[ui]}}
\index{configuration!ui@\texttt{[ui]}}\index{profile}%
One binary runs in two roles, selected by \cfg{[ui].profile} (overridable at
launch with \cfg{--profile}):
\begin{table}[h]
  \centering
  \begin{tabularx}{\textwidth}{KTD}
    \toprule
    \textbf{Key} & \textbf{Default} & \textbf{Meaning} \\
    \midrule
    profile & "pi" & \cfg{"pi"}: the deployed sorter --- an 800$\times$480
      touch panel with sorting controls only. \cfg{"workstation"}: a larger
      window that also exposes the learning interface and catalogue export,
      for teaching the system new objects. \\
    \bottomrule
  \end{tabularx}
  \caption{\cfg{[ui]} keys}
\end{table}

\section{Section \texttt{[recognition]}}
\index{configuration!recognition@\texttt{[recognition]}}\index{recognition}\index{object catalogue}%
Object recognition pairs an ONNX image embedder with a nearest-neighbour
\emph{catalogue} (the portable "learned file"): each object crop is embedded,
then matched by cosine similarity against the labelled examples you have
taught. It is \emph{disabled by default} --- nothing is recognised, and (since
sorting keys on class) nothing is picked, until you generate a model and teach
some classes. Recognition is independent of bins: see \S\ref{sec:calibration}
and the \cfg{[sorting]} assignments for how a recognised class is routed.
\begin{table}[h]
  \centering
  \begin{tabularx}{\textwidth}{KTD}
    \toprule
    \textbf{Key} & \textbf{Default} & \textbf{Meaning} \\
    \midrule
    enabled & false & Master switch. While false the classifier is a stub:
      every object is unrecognised. \\
    model\_path & "models/embedder.onnx" & ONNX embedder produced by
      \cfg{python models/export\_embedder.py} (see \cfg{models/README.md}). \\
    catalog\_path & "catalog.toml" & The learned catalogue. Portable --- copy
      it from the teaching workstation to the Pi. Absent = start empty. \\
    match\_threshold & 0.7 & Minimum cosine similarity (in $[-1,1]$) to accept
      a match; below it the object is left unrecognised. \\
    input\_size & 224 & Square input size the embedder expects; must match the
      export script. \\
    \bottomrule
  \end{tabularx}
  \caption{\cfg{[recognition]} keys}
\end{table}

\section{Camera calibration}
\label{sec:calibration}
\index{calibration}\index{camera!calibration}\index{coordinate transform}%
The pixel-to-robot transform is affine: scale (mm per pixel), rotation
and offset. There is deliberately no built-in default: at startup the
parameters are constructed from the capture resolution actually granted
by the camera and the configured \cfg{vision.mm\_per\_px} (the helper
\cfg{CalibrationParams::centered(width, height, mm\_per\_px)} maps the
optical centre to the robot origin, with rotation 0, i.e.\ the camera is
assumed to be mounted with its image y axis along the belt's $+Y$
direction). A guided calibration wizard that also measures rotation and
offset is planned (see the TODO list); until it exists, mount the camera
accordingly and measure \cfg{mm\_per\_px} by hand.

\section{Complete example}
\begin{lstlisting}[language=TOML]
[robot]
port_name = "/dev/ttyUSB0"
baud_rate = 115200
home_on_connect = true
x_min = -150.0
x_max = 150.0
y_min = -150.0
y_max = 150.0
z_min = -200.0
z_max = 0.0
z_pick = -180.0
z_travel = -150.0
feed_rate = 15000
release_gripper_on_estop = false

[conveyor]
port_name = "/dev/ttyUSB1"
baud_rate = 115200
default_speed = 1000
speed_mm_s = 100.0

[camera]
device_id = 0
width = 1280
height = 720

[[sorting.bins]]
id = "bin-1"
x = 120.0
y = 60.0
z = -100.0

[[sorting.assignments]]
class = "brick-2x4-red"
bin = "bin-1"

[vision]
threshold = 60.0
min_area = 500.0
max_area = 10000.0
invert = true
mm_per_px = 0.5

[logging]
level = "info"
directory = "logs"
to_console = true
keep_days = 7
\end{lstlisting}
